%% GENERATED by tools/gen_latex.py -- edit content/agents02.py instead.
%% Build:  latexmk -pdf agents02.tex     (or: pdflatex agents02.tex, twice)
\documentclass[aspectratio=169,11pt,t]{beamer}
\usepackage{itbpro}

\coursetitle{Model bahasa, dari sisi pembangun agen}
\coursesubtitle{Bukan cara kerja Transformer --- itu ada di Bab 15. Ini bagian yang mengubah keputusan: token sebagai satuan tagihan, konteks yang dibayar ulang, dan panggilan alat sebagai perilaku terlatih.}
\coursekicker{Chapter 2}
\coursesource{Grootendorst \& Alammar, An Illustrated Guide to AI Agents --- bab 2}
\courseduration{3 jam (2 sesi)}
\coursepresenter{Hendri Karisma}{Instructor}
\coursebrand{AI for Professional \textperiodcentered\ ITB \textperiodcentered\ Ch 2}

\title{Model bahasa, dari sisi pembangun agen}
\author{Hendri Karisma}
\date{}

\begin{document}
\covertitle

\framekicker{Learning outcomes}
\begin{frame}[shrink=25]{Session Objectives}
\leadpar{By the end of this session, participants can:}
\begin{isteps}
\item \textbf{Menghitung tagihan sebuah percakapan} dari panjang giliran dan jumlah giliran, dan menjelaskan kenapa ia tumbuh kuadrat.
\item \textbf{Menyebutkan apa yang dipertahankan singgahan prompt} dan satu kebiasaan yang mematikannya tanpa peringatan.
\item \textbf{Menjelaskan bahwa model tidak punya ingatan antar giliran}, dan menunjukkan di mana keadaan percakapan sebenarnya disimpan.
\item \textbf{Menerangkan panggilan alat sebagai keluaran terstruktur} yang belum dijalankan, dan menunjuk baris tempat ia dijalankan.
\item \textbf{Membedakan tiga tahap pelatihan} dan menyebutkan tahap mana yang membuat sebuah model bisa dipakai sebagai agen.
\item \textbf{Menyebutkan tiga sifat model} yang harus diukur sebelum memilihnya untuk sebuah agen --- dan mana yang tidak bisa dibaca dari papan skor.
\end{isteps}

\vskip 6pt
\vskip 4pt
\reslink{Site}{Course home}{https://hendrikarisma.my.id/rs/ai-products-course/index.html}
\reslink{Book}{Maarten Grootendorst \& Jay Alammar, An Illustrated Guide to AI Agents (O'Reilly Media, Early Release)}{https://www.oreilly.com/library/view/an-illustrated-guide/9798341662681/}
\vskip 2pt
\end{frame}

\framekicker{Bab ini}
\begin{frame}[shrink=25]{Yang sengaja TIDAK diulang di sini}
{\fontsize{9.4}{12.6}\selectfont\color{fg2}Kelas ini sudah punya \textbf{104 slide} tentang Transformer, perhatian diri, dan penyandian posisi di \textbf{Bab 15}, plus modul LLM tersendiri. Bab ini tidak mengulangnya.\par}\vskip 4pt
\begin{minipage}[t]{0.485\textwidth}
{\fontsize{9.4}{12.6}\selectfont\color{fg2}\textbf{Sudah dibahas di tempat lain}\par}\vskip 4pt
\begin{ibullets}
\item Cara perhatian diri menghitung --- Bab 15, dengan angkanya
\item Penyandian posisi, encoder lawan decoder --- Bab 15
\item Pelatihan awal dan penyetelan halus --- Bab 15--16
\end{ibullets}

\end{minipage}%
\hfill
\begin{minipage}[t]{0.485\textwidth}
{\fontsize{9.4}{12.6}\selectfont\color{fg2}\textbf{Yang dibahas di sini}\par}\vskip 4pt
\begin{ibullets}
\item Token sebagai satuan tagihan, waktu, dan batas
\item Kenapa percakapan panjang mahal secara kuadrat
\item Panggilan alat sebagai keluaran terstruktur
\item Sifat model mana yang menentukan keberhasilan agen
\end{ibullets}

\end{minipage}%
\begin{iband}[signal]
Aturannya sederhana: yang masuk ke sini hanya yang \textbf{mengubah keputusan} saat membangun agen.
\end{iband}
\end{frame}

\coursesection{01}{Token: satuan tagihan, waktu, dan batas}{Tiga hal yang berbeda, semuanya diukur dengan satuan yang sama.}

\framekicker{Token}
\begin{frame}[shrink=25]{Model tidak melihat kata, dan tidak melihat huruf}
{\fontsize{9.4}{12.6}\selectfont\color{fg2}Teks dipecah jadi \textbf{token} --- potongan yang lebih besar dari huruf dan biasanya lebih kecil dari kata. Kata yang sering muncul jadi satu token; kata yang jarang pecah jadi beberapa.\par}\vskip 4pt
{\fontsize{9.4}{12.6}\selectfont\color{fg2}Kenapa ini penting bagi pembangun agen, bukan bagi ahli bahasa: \textbf{semua yang Anda bayar, semua yang Anda tunggu, dan semua batas yang Anda tabrak dihitung dalam token} --- bukan dalam karakter, kalimat, atau permintaan.\par}\vskip 4pt
\begin{iband}[signal]
Taksiran kasar yang cukup untuk perencanaan: \textbf{satu token $\approx$ 4 karakter} untuk teks Inggris, dan lebih boros untuk bahasa Indonesia serta untuk JSON --- dan hasil alat Anda hampir selalu JSON.
\end{iband}
\end{frame}

\framekicker{Token}
\begin{frame}[shrink=25]{Masukan dan keluaran tidak berharga sama, dan tidak berperilaku sama}
\begin{minipage}[t]{0.485\textwidth}
{\fontsize{9.4}{12.6}\selectfont\color{fg2}\textbf{Token masukan} --- perintah sistem, skema alat, seluruh riwayat, semua hasil alat. Dibaca sekaligus, sejajar, jadi cepat.\par}\vskip 4pt
{\fontsize{9.4}{12.6}\selectfont\color{fg2}Di sebuah agen, inilah yang membengkak. Bukan jawabannya.\par}\vskip 4pt

\end{minipage}%
\hfill
\begin{minipage}[t]{0.485\textwidth}
{\fontsize{9.4}{12.6}\selectfont\color{fg2}\textbf{Token keluaran} --- yang ditulis model. Dihasilkan satu per satu, tiap token menunggu token sebelumnya, jadi \textbf{lambat}.\par}\vskip 4pt
{\fontsize{9.4}{12.6}\selectfont\color{fg2}Lebih mahal per token, tapi di sebuah agen jumlahnya jauh lebih sedikit: panggilan alat itu pendek.\par}\vskip 4pt

\end{minipage}%
{\fontsize{9.4}{12.6}\selectfont\color{fg2}Akibat praktisnya: \textbf{agen adalah beban masukan.} Mengoptimalkan panjang jawaban hampir tidak mengubah apa pun; memendekkan yang masuk mengubah semuanya.\par}\vskip 4pt
\end{frame}

\framekicker{Token}
\begin{frame}[shrink=25]{Percakapannya tumbuh lurus, tagihannya tidak}
\begin{center}
\includegraphics[width=0.97\linewidth,height=0.56\textheight,keepaspectratio]{figs/agents02-context.tex.pdf}
\end{center}
\vskip 3pt{\color{ink3}\fontsize{7.4}{9.4}\selectfont Sepuluh giliran, dihitung dari perintah sistem 800 token dan 350 token per giliran. Langkahi: giliran pertama, lima pertama, lalu sepuluh.\par}
{\fontsize{9.4}{12.6}\selectfont\color{fg2}Percakapan yang berakhir pada 3 950 token menagihkan \textbf{23 750} token masukan --- enam kali lipat. Dan ini bukan sifat khusus agen; ini sifat percakapan mana pun. Yang khusus pada agen adalah \textbf{giliran per tugasnya banyak}, jadi ia berada jauh di kanan grafik itu.\par}\vskip 4pt
\end{frame}

\framekicker{Token}
\begin{frame}[shrink=25]{Dua obatnya, dan yang kedua sering dirusak sendiri}
\begin{minipage}[t]{0.494\textwidth}
\begin{icardaccent}
\cardh{Kurangi gilirannya}
Satu langkah yang dihapus menghemat lebih banyak daripada seribu token yang dipangkas, sebab langkah itu dibayar ulang oleh \textbf{semua} langkah sesudahnya.
\end{icardaccent}
\end{minipage}%
\hfill
\begin{minipage}[t]{0.494\textwidth}
\begin{icardaccent}
\cardh{Singgahkan bagian yang tetap}
Perintah sistem, skema alat, dan teks kebijakan tidak berubah antar giliran. Penyedia menagihnya jauh lebih murah kalau awalannya \textbf{sama persis}.
\end{icardaccent}
\end{minipage}%
\begin{iband}[signal]
Satu stempel waktu di dalam perintah sistem membuat awalannya berbeda tiap kali, dan \hilite{singgahannya mati tanpa satu pun pesan galat}. Periksa \inlinecode{cache\_read\_input\_tokens} bukan nol --- jangan percaya bahwa ia menyala.
\end{iband}
\end{frame}

\framekicker{Token}
\begin{frame}[shrink=25]{Jendela konteks adalah dinding, bukan saran}
{\fontsize{9.4}{12.6}\selectfont\color{fg2}Tiap model punya batas jumlah token yang boleh masuk sekaligus. Melewatinya bukan menghasilkan jawaban yang lebih buruk --- ia menghasilkan \textbf{galat}, atau diam-diam memotong bagian paling awal percakapan.\par}\vskip 4pt
{\fontsize{9.4}{12.6}\selectfont\color{fg2}Yang paling awal biasanya \textbf{perintah sistem dan batasan Anda}. Jadi kegagalannya berbentuk: agen bekerja normal selama sepuluh giliran, lalu pada giliran ke dua puluh mulai melanggar aturan yang sudah tidak ada lagi di konteksnya.\par}\vskip 4pt
\begin{iband}[signal]
Karena itu memangkas riwayat harus jadi \textbf{keputusan yang Anda tulis}, bukan akibat sampingan. Bab 4 membahas apa yang disimpan dan apa yang dibuang.
\end{iband}
\end{frame}

\coursesection{02}{Dari model bahasa ke mesin agen}{Tiga hal yang harus ada sebelum sebuah model bisa dipakai dalam gelung.}

\framekicker{Dari model ke agen}
\begin{frame}[shrink=25]{Model tidak punya ingatan --- riwayatnya yang punya}
\begin{center}
\includegraphics[width=0.92\linewidth,height=0.60\textheight,keepaspectratio]{figs/agents02-stateless.tex.pdf}
\end{center}
\vskip 3pt{\color{ink3}\fontsize{7.4}{9.4}\selectfont Dua giliran adalah dua permintaan yang sama sekali terpisah.\par}
{\fontsize{9.4}{12.6}\selectfont\color{fg2}Antara dua giliran, model \textbf{tidak menyimpan apa pun}. Yang membuatnya tampak mengingat adalah kode Anda mengirim ulang seluruh percakapan tiap kali. Keadaan percakapan ada di \textbf{daftar pesan milik Anda}, dan itu kabar baik: keadaan yang Anda pegang adalah keadaan yang bisa Anda periksa, potong, dan simpan.\par}\vskip 4pt
\end{frame}

\framekicker{Dari model ke agen}
\begin{frame}[shrink=25]{Perintah sistem: kontrak, bukan mantra}
{\fontsize{9.4}{12.6}\selectfont\color{fg2}Perintah sistem menetapkan peran, gaya, dan batasan. Ia berguna dan perlu ditulis dengan hati-hati --- tetapi ia \textbf{bukan mekanisme keamanan}.\par}\vskip 4pt
\begin{minipage}[t]{0.485\textwidth}
{\fontsize{9.4}{12.6}\selectfont\color{fg2}\textbf{Yang pantas ditaruh di sana}\par}\vskip 4pt
\begin{ibullets}
\item Peran dan pembacanya
\item Format jawaban yang diharapkan
\item Kapan harus menyerah dan bertanya
\item Kebijakan yang harus dikutip, bukan diingat
\end{ibullets}

\end{minipage}%
\hfill
\begin{minipage}[t]{0.485\textwidth}
{\fontsize{9.4}{12.6}\selectfont\color{fg2}\textbf{Yang tidak boleh hanya di sana}\par}\vskip 4pt
\begin{ibullets}
\item ``Jangan pernah menghapus data'' $\rightarrow$ jangan beri alat hapus
\item ``Jangan bocorkan data pribadi'' $\rightarrow$ jangan beri alat yang bisa membacanya
\item ``Jangan menyetujui'' $\rightarrow$ jangan sediakan alat menyetujui
\end{ibullets}

\end{minipage}%
\begin{iband}[signal]
Uji sederhana: kalau sebuah kalimat di perintah sistem hilang, apakah sistemnya jadi tidak aman? Kalau ya, \hilite{kalimat itu berada di tempat yang salah}.
\end{iband}
\end{frame}

\framekicker{Dari model ke agen}
\begin{frame}[shrink=25]{Panggilan alat adalah keluaran terstruktur, bukan eksekusi}
\begin{center}
\includegraphics[width=0.92\linewidth,height=0.60\textheight,keepaspectratio]{figs/agents02-toolcall.tex.pdf}
\end{center}
\vskip 3pt{\color{ink3}\fontsize{7.4}{9.4}\selectfont Model menghasilkan niat; kode memutuskan apakah niat itu dijalankan.\par}
{\fontsize{9.4}{12.6}\selectfont\color{fg2}Model diberi \textbf{skema} tiap alat --- namanya, kegunaannya, parameternya beserta tipe. Yang ia keluarkan hanyalah sebuah objek: nama alat dan argumen. Objek itu belum melakukan apa-apa.\par}\vskip 4pt
\begin{iband}[signal]
Ini fakta arsitektural terpenting dalam seluruh modul, dan ia muncul lagi di tiap bab: \textbf{model menuliskan niat, kode Anda yang mengeksekusi.} Semua batas izin yang Anda punya berada di celah itu.
\end{iband}
\end{frame}

\framekicker{Dari model ke agen}
\begin{frame}[shrink=25]{Deskripsi alat adalah antarmuka, dan ditulis untuk pembaca yang aneh}
{\fontsize{9.4}{12.6}\selectfont\color{fg2}Satu-satunya yang dipakai model untuk memilih alat adalah \textbf{deskripsinya}. Bukan kodenya, bukan namanya saja. Deskripsi yang buruk menghasilkan agen yang tampak bodoh, dan orang biasanya menyalahkan modelnya.\par}\vskip 4pt
\begin{itable}{>{\raggedright\arraybackslash}p{0.282\textwidth}>{\raggedright\arraybackslash}p{0.3384\textwidth}>{\raggedright\arraybackslash}p{0.3196\textwidth}}
\thead{Ditulis begini} & \thead{Akibatnya} & \thead{Perbaikannya} \\ \midrule
\inlinecode{cari(q)} --- mencari data & Dipanggil untuk apa saja, termasuk yang bukan urusannya & Sebut \textbf{apa} yang dicari dan \textbf{kapan} memakainya \\
Dua alat dengan deskripsi mirip & Dipilih acak antara keduanya & Gabungkan, atau bedakan dengan tegas \\
Parameter tanpa tipe atau contoh & Argumen ngawur, galat di dalam alat & Tipe ketat + satu contoh nilai \\
Deskripsi menyebut cara kerjanya & Tidak membantu memilih & Sebut \textbf{hasilnya}, bukan implementasinya \\
\end{itable}
\end{frame}

\framekicker{Dari model ke agen}
\begin{frame}[shrink=25]{Suhu: satu angka yang diam-diam menentukan keterulangan}
{\fontsize{9.4}{12.6}\selectfont\color{fg2}Model memilih token berikutnya dari sebaran peluang. \textbf{Suhu} mengatur seberapa berani pilihannya menyimpang dari yang paling mungkin.\par}\vskip 4pt
\begin{minipage}[t]{0.3253\textwidth}
\begin{icardgood}
\cardh{Rendah (0--0,3)}
Nyaris sama tiap kali. Yang Anda mau untuk \textbf{memilih alat} dan mengeluarkan JSON.
\end{icardgood}
\end{minipage}%
\hfill
\begin{minipage}[t]{0.3253\textwidth}
\begin{icard}
\cardh{Sedang (0,5--0,7)}
Untuk menyusun penjelasan yang enak dibaca, sesudah keputusannya diambil.
\end{icard}
\end{minipage}%
\hfill
\begin{minipage}[t]{0.3253\textwidth}
\begin{icardbad}
\cardh{Tinggi (> 1)}
Untuk mencari ragam gagasan. Hampir tidak pernah untuk agen yang menyentuh sistem nyata.
\end{icardbad}
\end{minipage}%
\begin{iband}[signal]
Dan satu jebakan pengujian: pada suhu tinggi, \textbf{kumpulan uji yang lulus hari ini bisa gagal besok tanpa satu baris pun berubah}. Kalau hasil uji Anda goyah, periksa angka ini sebelum menyalahkan agennya.
\end{iband}
\end{frame}

\framekicker{Token}
\begin{frame}[shrink=25]{Hasil alat adalah JSON, dan JSON itu boros}
{\fontsize{9.4}{12.6}\selectfont\color{fg2}Teks prosa mendekati 4 karakter per token. JSON tidak: tanda kutip, kurung kurawal, koma, dan nama kunci yang berulang di tiap baris semuanya dihitung.\par}\vskip 4pt
{\fontsize{9.4}{12.6}\selectfont\color{fg2}Satu tabel 20 baris dengan 8 kolom bisa menghabiskan beberapa ribu token --- dan karena ia masuk riwayat, ia \textbf{dibayar ulang di tiap giliran sesudahnya}. Ini penyebab pembengkakan biaya yang paling sering tidak terlihat, sebab yang dilihat orang adalah panjang jawabannya.\par}\vskip 4pt
\begin{minipage}[t]{0.3253\textwidth}
\begin{icardgood}
\cardh{Kembalikan yang dipakai}
Alat yang mengembalikan seluruh baris ketika agen hanya butuh satu angka membayar seribu token untuk satu.
\end{icardgood}
\end{minipage}%
\hfill
\begin{minipage}[t]{0.3253\textwidth}
\begin{icardgood}
\cardh{Ringkas di sisi alat}
Peringkasan yang dilakukan \textbf{kode} gratis dan pasti; yang dilakukan model berbayar dan bisa salah.
\end{icardgood}
\end{minipage}%
\hfill
\begin{minipage}[t]{0.3253\textwidth}
\begin{icard}
\cardh{Rujuk, jangan salin}
Simpan hasil besar di luar, kembalikan pengenalnya. Bab 4 menyebut ini memori kerja.
\end{icard}
\end{minipage}%
\end{frame}

\framekicker{Token}
\begin{frame}[shrink=25]{Jendela konteks itu anggaran, dan sebagian sudah habis sebelum mulai}
\begin{center}
\includegraphics[width=0.97\linewidth,height=0.56\textheight,keepaspectratio]{figs/agents02-budget.tex.pdf}
\end{center}
\vskip 3pt{\color{ink3}\fontsize{7.4}{9.4}\selectfont Jendela 128k, dibelanjakan: bagian tetap, lalu hasil alat yang menumpuk tiap giliran. Langkahi 5, 15, dan 30 giliran.\par}
{\fontsize{9.4}{12.6}\selectfont\color{fg2}Sepuluh ribu token sudah terpakai \textbf{sebelum giliran pertama}, dan tiap alat baru memotong lagi. Itu sebabnya ``tambahkan saja satu alat lagi'' bukan keputusan gratis: ia memperkecil ruang kerja tiap giliran sesudahnya.\par}\vskip 4pt
\end{frame}

\framekicker{Token}
\begin{frame}[shrink=25]{Satu giliran agen adalah dua perjalanan penuh}
\begin{center}
\includegraphics[width=0.92\linewidth,height=0.60\textheight,keepaspectratio]{figs/agents02-roundtrip.tex.pdf}
\end{center}
\vskip 3pt{\color{ink3}\fontsize{7.4}{9.4}\selectfont Memanggil satu alat berarti mengirim seluruh riwayat dua kali.\par}
{\fontsize{9.4}{12.6}\selectfont\color{fg2}Model memutuskan alat mana yang dipanggil --- itu satu permintaan. Hasilnya dibaca dan diubah jadi jawaban --- itu permintaan kedua, dengan riwayat yang sekarang lebih panjang. \textbf{Satu panggilan alat berharga dua permintaan}, bukan satu.\par}\vskip 4pt
\begin{iband}[signal]
Karena itu ``berapa alat yang dipanggil'' adalah ukuran biaya yang lebih baik daripada ``berapa permintaan'', dan kenapa menggabungkan dua alat kecil jadi satu sering lebih hemat daripada mengoptimalkan keduanya.
\end{iband}
\end{frame}

\framekicker{Waktu}
\begin{frame}[shrink=25]{Waktu tunggu bukan satu angka}
\begin{minipage}[t]{0.485\textwidth}
{\fontsize{9.4}{12.6}\selectfont\color{fg2}\textbf{Waktu ke token pertama} --- dipengaruhi panjang masukan. Konteks panjang membuat jeda pertama terasa, dan inilah yang dirasakan pengguna sebagai ``lambat''.\par}\vskip 4pt

\end{minipage}%
\hfill
\begin{minipage}[t]{0.485\textwidth}
{\fontsize{9.4}{12.6}\selectfont\color{fg2}\textbf{Waktu total} --- dipengaruhi panjang keluaran, sebab token dihasilkan satu per satu. Di agen ini biasanya kecil; yang besar adalah \textbf{waktu alat}.\par}\vskip 4pt

\end{minipage}%
{\fontsize{9.4}{12.6}\selectfont\color{fg2}Pada agen enam langkah, sebagian besar waktu dinding sering bukan milik model sama sekali --- melainkan milik kueri basis data dan panggilan HTTP di dalam alat. Mengukur waktu model saja akan mengoptimalkan bagian yang salah.\par}\vskip 4pt
\begin{iband}[signal]
Ukur per \textbf{giliran}, bukan per tugas: waktu model, waktu alat, dan jumlah giliran. Tiga angka itu memberi tahu di mana waktunya pergi; satu angka total tidak.
\end{iband}
\end{frame}

\framekicker{Waktu}
\begin{frame}[shrink=25]{Mengalirkan jawaban tidak membuat agen lebih cepat}
{\fontsize{9.4}{12.6}\selectfont\color{fg2}Pengaliran (\textit{streaming}) mengirim token begitu dihasilkan. Untuk obrolan itu besar pengaruhnya pada rasa: pengguna melihat sesuatu bergerak.\par}\vskip 4pt
{\fontsize{9.4}{12.6}\selectfont\color{fg2}Untuk agen, pengaruhnya jauh lebih kecil, dan sering \textbf{nol} --- sebab yang keluar dari model bukan jawaban, melainkan panggilan alat yang harus \textbf{lengkap} sebelum bisa dijalankan. Tidak ada yang bisa dikerjakan dari separuh panggilan alat.\par}\vskip 4pt
\begin{iband}[signal]
Yang benar-benar mengurangi rasa lambat pada agen: \textbf{menampilkan alat mana yang sedang berjalan}. Itu sebabnya layar ketiga di demo menyebut nama alatnya, bukan memutar lingkaran.
\end{iband}
\end{frame}

\framekicker{Keluaran terstruktur}
\begin{frame}[shrink=25]{Tiga cara memaksa bentuk, dengan tiga tingkat jaminan}
\begin{itable}{>{\raggedright\arraybackslash}p{0.2444\textwidth}>{\raggedright\arraybackslash}p{0.376\textwidth}>{\raggedright\arraybackslash}p{0.3196\textwidth}}
\thead{Cara} & \thead{Jaminannya} & \thead{Kapan dipakai} \\ \midrule
Diminta di prompt & Tidak ada. Sering benar, kadang tidak, dan gagalnya tidak seragam & Prototipe, atau model tanpa dukungan lain \\
Mode JSON penyedia & JSON-nya sah, \textbf{isinya belum tentu sesuai skema} & Ketika hanya butuh objek yang bisa diurai \\
Panggilan alat / skema ketat & Nama dan tipe parameter dijamin di sisi penyedia & Untuk agen --- ini yang dipakai \\
\end{itable}
{\fontsize{9.4}{12.6}\selectfont\color{fg2}Perbedaan baris kedua dan ketiga adalah perbedaan antara \textit{bisa diurai} dan \textit{bisa dipercaya}. Yang pertama menghindarkan galat penguraian; hanya yang kedua yang menghindarkan argumen yang masuk akal secara sintaks dan salah secara isi.\par}\vskip 4pt
\end{frame}

\framekicker{Keluaran terstruktur}
\begin{frame}[shrink=25]{Validasi tetap di sisi Anda, dan gagalnya harus punya jalan}
{\fontsize{9.4}{12.6}\selectfont\color{fg2}Skema ketat menjamin \textbf{bentuk}, bukan \textbf{kebenaran}. \inlinecode{id\_pelanggan} bertipe string yang berisi id milik orang lain tetap lolos skema.\par}\vskip 4pt
\begin{isteps}
\item \{'h': 'Periksa di batas alat, sebelum efek apa pun', 'p': 'Rentang, keberadaan, kepemilikan. Alat yang dipanggil dengan argumen tidak sah harus \textbf{menolak}, bukan menebak maksudnya.'\}
\item \{'h': 'Kembalikan galat yang bisa dipakai model', 'p': '``id tidak ditemukan'' membuatnya mencoba jalan lain. ``Galat 500'' membuatnya mengulang hal yang sama.'\}
\item \{'h': 'Hitung berapa kali ini terjadi', 'p': 'Laju penolakan skema adalah ukuran kualitas model yang paling cepat memberi tahu Anda telah memilih model yang salah.'\}
\end{isteps}
\end{frame}

\framekicker{Keluaran terstruktur}
\begin{frame}[shrink=25]{Determinisme yang cukup untuk bisa diuji}
{\fontsize{9.4}{12.6}\selectfont\color{fg2}Kumpulan uji yang hasilnya berubah-ubah tanpa perubahan kode bukan kumpulan uji --- ia sumber kebisingan yang membuat orang berhenti mempercayainya.\par}\vskip 4pt
\begin{minipage}[t]{0.494\textwidth}
\begin{icardgood}
\cardh{Suhu nol untuk jalur keputusan}
Pemilihan alat dan pengeluaran argumen tidak butuh keragaman. Keragaman di sana namanya ketidakstabilan.
\end{icardgood}
\end{minipage}%
\hfill
\begin{minipage}[t]{0.494\textwidth}
\begin{icardgood}
\cardh{Penyedia luring untuk menguji mesinnya}
Gelung, anggaran, penjaga, dan validasi alat tidak butuh model sungguhan untuk diuji --- dan uji yang tidak butuh jaringan akan benar-benar dijalankan orang.
\end{icardgood}
\end{minipage}%
\begin{iband}[signal]
Bahkan pada suhu nol, penyedia \textbf{tidak menjanjikan} keluaran yang identik antar waktu. Uji yang mensyaratkan kalimat persis akan rapuh; uji yang mensyaratkan \textbf{alat mana yang dipanggil} tidak.
\end{iband}
\end{frame}

\coursesection{03}{Bagaimana ia jadi bisa disuruh}{Tiga tahap, dan hanya dua terakhir yang membuatnya berguna sebagai agen.}

\framekicker{Pelatihan}
\begin{frame}[shrink=25]{Tiga tahap, tiga kemampuan yang berbeda}
\begin{center}
\includegraphics[width=0.92\linewidth,height=0.60\textheight,keepaspectratio]{figs/agents02-training.tex.pdf}
\end{center}
\vskip 3pt{\color{ink3}\fontsize{7.4}{9.4}\selectfont Tiap tahap menambah kemampuan yang tidak dimiliki tahap sebelumnya.\par}
{\fontsize{9.4}{12.6}\selectfont\color{fg2}\textbf{Pra-latih} memberinya bahasa dan pengetahuan dunia, dari meramal token berikutnya pada teks dalam jumlah raksasa. Model di tahap ini belum bisa disuruh --- ia melanjutkan teks, tidak menjawab pertanyaan.\par}\vskip 4pt
\end{frame}

\framekicker{Pelatihan}
\begin{frame}[shrink=25]{Tahap yang membuatnya bisa dipakai dalam gelung}
\begin{minipage}[t]{0.485\textwidth}
{\fontsize{9.4}{12.6}\selectfont\color{fg2}\textbf{Penyetelan terbimbing (SFT)}\par}\vskip 4pt
{\fontsize{9.4}{12.6}\selectfont\color{fg2}Dilatih pada contoh percakapan yang ditulis manusia: permintaan, dan jawaban yang pantas. Di sinilah ia belajar bentuk \textit{menjawab}, dan di sini pula ia belajar bentuk \textbf{panggilan alat} --- sebab contohnya berisi panggilan alat.\par}\vskip 4pt

\end{minipage}%
\hfill
\begin{minipage}[t]{0.485\textwidth}
{\fontsize{9.4}{12.6}\selectfont\color{fg2}\textbf{Penguatan dari umpan balik}\par}\vskip 4pt
{\fontsize{9.4}{12.6}\selectfont\color{fg2}Manusia (atau model penilai) membandingkan dua jawaban; model diberi imbalan atas yang lebih disukai. Di sinilah ia belajar \textbf{menuruti batasan}, menolak, dan berhenti mengarang dengan percaya diri.\par}\vskip 4pt

\end{minipage}%
\begin{iband}[signal]
Kenapa ini penting bagi pembangun agen: \textbf{kepatuhan pada skema alat dan pada instruksi adalah hasil pelatihan, bukan sifat bawaan Transformer.} Dua model dengan arsitektur nyaris sama bisa sangat berbeda di sini, dan itulah yang paling menentukan.
\end{iband}
\end{frame}

\framekicker{Pelatihan}
\begin{frame}[shrink=25]{Yang tidak diperbaiki oleh tahap mana pun}
{\fontsize{9.4}{12.6}\selectfont\color{fg2}Tidak satu pun dari ketiga tahap itu membuat model \textbf{tahu apa yang tidak diketahuinya}. Ia dilatih menghasilkan lanjutan yang masuk akal, dan lanjutan yang masuk akal tetap masuk akal ketika salah.\par}\vskip 4pt
\begin{isteps}
\item \{'h': 'Akibat pertama: percaya diri bukan sinyal', 'p': 'Nada yakin dihasilkan oleh pelatihan yang sama, entah isinya benar atau tidak.'\}
\item \{'h': 'Akibat kedua: sumber harus dari alat', 'p': 'Angka yang tidak berasal dari hasil alat adalah angka yang tidak punya asal-usul. Itu sebabnya bab 5 memperlakukan alat sebagai sumber bukti, bukan sekadar kemampuan.'\}
\item \{'h': 'Akibat ketiga: penilaian harus melihat jejak', 'p': 'Keluaran yang benar dari alasan yang salah akan berulang. Bab 7.'\}
\end{isteps}
\end{frame}

\coursesection{04}{Arsitektur, sejauh yang mengubah keputusan}{Kedalamannya di Bab 15. Di sini hanya bagian yang terasa di tagihan.}

\framekicker{Arsitektur}
\begin{frame}[shrink=25]{Cache KV: kenapa token pertama lambat dan sisanya cepat}
{\fontsize{9.4}{12.6}\selectfont\color{fg2}Untuk menghasilkan token, model melihat seluruh token sebelumnya. Menghitung ulang semuanya tiap token akan sangat boros, jadi hasil antaranya \textbf{disimpan} --- itulah cache KV.\par}\vskip 4pt
\begin{minipage}[t]{0.485\textwidth}
{\fontsize{9.4}{12.6}\selectfont\color{fg2}\textbf{Fase isi} --- seluruh masukan diproses sekaligus. Sejajar, cepat per token, tapi inilah yang membuat token pertama terasa lama pada konteks panjang.\par}\vskip 4pt

\end{minipage}%
\hfill
\begin{minipage}[t]{0.485\textwidth}
{\fontsize{9.4}{12.6}\selectfont\color{fg2}\textbf{Fase lanjut} --- satu token per langkah, memakai cache. Murah dalam komputasi, tetapi cachenya \textbf{memakan memori} yang tumbuh bersama panjang konteks.\par}\vskip 4pt

\end{minipage}%
\begin{iband}[signal]
Yang terasa di produksi: konteks panjang bukan hanya lebih mahal ditagih --- ia juga \textbf{menurunkan jumlah permintaan bersamaan} yang muat di satu GPU. Itu sebabnya penyedia membedakan harga menurut panjang konteks.
\end{iband}
\end{frame}

\framekicker{Arsitektur}
\begin{frame}[shrink=25]{Campuran pakar: kenapa model besar bisa murah}
{\fontsize{9.4}{12.6}\selectfont\color{fg2}Sebagian model modern hanya menyalakan \textbf{sebagian kecil} parameternya untuk tiap token. Jumlah parameter totalnya besar; yang aktif per token jauh lebih kecil.\par}\vskip 4pt
{\fontsize{9.4}{12.6}\selectfont\color{fg2}Akibatnya satu angka yang sering dipakai membandingkan model --- jumlah parameter --- \textbf{berhenti berarti apa-apa} sebagai penanda biaya atau kecepatan. Dua model dengan angka yang sama bisa berbeda berkali lipat ongkosnya.\par}\vskip 4pt
\begin{iband}[signal]
Aturan praktisnya: bandingkan model dengan \textbf{biaya per tugas yang selesai}, diukur pada kumpulan uji Anda sendiri. Bukan dengan jumlah parameter, dan bukan dengan harga per juta token.
\end{iband}
\end{frame}

\framekicker{Arsitektur}
\begin{frame}[shrink=25]{Ke mana mencari kedalamannya}
\begin{itable}{>{\raggedright\arraybackslash}p{0.3196\textwidth}>{\raggedright\arraybackslash}p{0.2068\textwidth}>{\raggedright\arraybackslash}p{0.4136\textwidth}}
\thead{Kalau ingin tahu} & \thead{Lihat} & \thead{Yang ada di sana} \\ \midrule
Cara perhatian diri menghitung & \textbf{Bab 15} & Q$\cdot$K$\cdot$V pada kalimat nyata sampai persentase softmax, dan satu slide yang membaca angkanya dengan jujur \\
Kenapa perlu proyeksi Q dan K & \textbf{Bab 15} & Jejak \inlinecode{run}: tanpa proyeksi, sebuah kata menghabiskan 60\% perhatiannya untuk dirinya sendiri \\
Penyandian posisi & \textbf{Bab 15} & Gambar yang menghitung sendiri \\
Menghasilkan teks, pengambilan sampel & \textbf{Bab 16} & Suhu, top-p, dan gelung pengambilan sampel \\
\end{itable}
{\fontsize{9.4}{12.6}\selectfont\color{fg2}Dek ini sengaja berhenti di sini. Menuliskan ulang isi Bab 15 dengan kata yang berbeda akan menambah slide tanpa menambah apa pun yang bisa dipakai.\par}\vskip 4pt
\end{frame}

\coursesection{05}{Memilih model untuk sebuah agen}{Papan skor menjawab pertanyaan yang berbeda dari pertanyaan Anda.}

\framekicker{Memilih}
\begin{frame}[shrink=25]{Tiga sifat yang menentukan, dan hanya satu yang ada di papan skor}
\begin{minipage}[t]{0.3253\textwidth}
\begin{icardaccent}
\cardh{Kepatuhan pada skema}
Seberapa sering ia mengeluarkan panggilan alat yang \textbf{sah} dan argumen yang bertipe benar. Kegagalan di sini terlihat seperti agen yang bodoh, padahal soal format.
\end{icardaccent}
\end{minipage}%
\hfill
\begin{minipage}[t]{0.3253\textwidth}
\begin{icardaccent}
\cardh{Kemauan berhenti}
Seberapa sering ia mengatakan \textit{tidak bisa} alih-alih menebak. Model yang tidak pernah menyerah menghasilkan agen yang berputar.
\end{icardaccent}
\end{minipage}%
\hfill
\begin{minipage}[t]{0.3253\textwidth}
\begin{icardaccent}
\cardh{Kepatuhan pada instruksi panjang}
Batasan di token ke-3 000 masih dituruti pada giliran ke-15? Inilah yang paling cepat rusak saat konteks memanjang.
\end{icardaccent}
\end{minipage}%
{\fontsize{9.4}{12.6}\selectfont\color{fg2}Papan skor umum mengukur pengetahuan dan penalaran --- berguna, tetapi bukan tiga hal di atas. Ketiganya harus \textbf{diukur pada kumpulan uji Anda sendiri}, dan itu satu sore kerja, bukan satu proyek.\par}\vskip 4pt
\end{frame}

\framekicker{Memilih}
\begin{frame}[shrink=25]{Urutan yang menghemat paling banyak waktu}
\begin{isteps}
\item \{'h': 'Mulai dari model paling mampu yang ada', 'p': 'Sampai agennya \textbf{benar}. Mengoptimalkan agen yang belum benar hanya membuat kesalahannya lebih murah.'\}
\item \{'h': 'Baru turunkan, satu tingkat, dengan kumpulan uji di tangan', 'p': 'Kalau angkanya bertahan, simpan penghematannya. Kalau tidak, Anda baru saja mengetahui harga sebenarnya dari model yang lebih murah.'\}
\item \{'h': 'Ukur biaya per tugas selesai, bukan per permintaan', 'p': 'Model murah yang perlu dua kali lebih banyak giliran tidak menghemat apa pun.'\}
\item \{'h': 'Pertimbangkan model lokal ketika datanya yang menentukan', 'p': 'Bukan karena lebih murah --- sering tidak. Karena ada data yang tidak boleh keluar, dan itu keputusan kepatuhan, bukan teknis.'\}
\end{isteps}
\begin{iband}[signal]
\inlinecode{ai-agentic-demo} menjalankan kelimanya lewat satu antarmuka: \inlinecode{echo} (luring), \inlinecode{ollama} (mesin sendiri), \inlinecode{gemini}, \inlinecode{anthropic}, \inlinecode{openai} --- jadi menukar model adalah satu argumen, bukan satu penulisan ulang.
\end{iband}
\end{frame}

\framekicker{Memilih}
\begin{frame}[shrink=25]{Bahasa Indonesia lebih mahal per kalimat, dan itu bukan dugaan}
{\fontsize{9.4}{12.6}\selectfont\color{fg2}Pemenggal token pada kebanyakan model dilatih pada teks yang didominasi bahasa Inggris. Kata Inggris yang umum sering jadi \textbf{satu} token; kata Indonesia yang sama umumnya pecah jadi dua, tiga, atau lebih --- terutama kata berimbuhan.\par}\vskip 4pt
{\fontsize{9.4}{12.6}\selectfont\color{fg2}Akibatnya langsung dan berlipat: kalimat yang sama artinya menghabiskan lebih banyak token, jadi lebih mahal, lebih cepat memenuhi jendela konteks, dan lebih lambat dihasilkan.\par}\vskip 4pt
\begin{iband}[signal]
Jangan menaksir --- \textbf{ukur pada teks Anda sendiri}. Ambil seratus kalimat nyata dari domain Anda, hitung tokennya dengan pemenggal model yang dipakai, dan pakai angka itu untuk perencanaan biaya. Taksiran 4 karakter per token adalah taksiran untuk bahasa Inggris.
\end{iband}
\end{frame}

\framekicker{Memilih}
\begin{frame}[shrink=25]{Jendela besar tidak berarti seluruhnya terpakai dengan baik}
{\fontsize{9.4}{12.6}\selectfont\color{fg2}Jendela satu juta token tidak berarti model memperhatikan satu juta token sama baiknya. Kemampuan menemukan informasi cenderung \textbf{paling baik di awal dan di akhir} konteks, dan paling lemah di tengah.\par}\vskip 4pt
\begin{minipage}[t]{0.485\textwidth}
{\fontsize{9.4}{12.6}\selectfont\color{fg2}\textbf{Akibat untuk penyusunan konteks}\par}\vskip 4pt
\begin{ibullets}
\item Taruh perintah dan batasan di \textbf{awal}
\item Taruh permintaan giliran ini di \textbf{akhir}
\item Jangan menaruh aturan penting di tengah tumpukan hasil alat
\end{ibullets}

\end{minipage}%
\hfill
\begin{minipage}[t]{0.485\textwidth}
{\fontsize{9.4}{12.6}\selectfont\color{fg2}\textbf{Akibat untuk arsitektur}\par}\vskip 4pt
\begin{ibullets}
\item Konteks besar bukan pengganti pengambilan yang tepat
\item Memasukkan seluruh dokumen karena ``muat'' menurunkan ketepatan, bukan menaikkannya
\item Bab 4 membahas apa yang disimpan dan apa yang dicari lagi
\end{ibullets}

\end{minipage}%
\begin{iband}[signal]
Uji ini murah dan jarang dilakukan: \hilite{sisipkan satu fakta di tengah konteks panjang dan tanyakan} --- pada panjang yang benar-benar Anda pakai, bukan pada panjang maksimum di brosur.
\end{iband}
\end{frame}

\framekicker{Operasional}
\begin{frame}[shrink=25]{Model diperbarui, dan agen Anda berubah tanpa Anda menyentuhnya}
{\fontsize{9.4}{12.6}\selectfont\color{fg2}Nama model tanpa versi menunjuk ke sasaran yang bergerak. Penyedia memperbarui bobotnya; perilaku pemilihan alat, panjang jawaban, dan kepatuhan format ikut bergeser --- tanpa satu baris pun berubah di sisi Anda.\par}\vskip 4pt
\begin{isteps}
\item \{'h': 'Sematkan versinya di produksi', 'p': 'Kalau penyedianya menyediakan pengenal versi, pakai. Kalau tidak, catat tanggal dan jalankan kumpulan uji lebih sering.'\}
\item \{'h': 'Catat versi model pada tiap jejak', 'p': 'Tanpa itu, ``sejak kapan ini mulai salah'' tidak bisa dijawab. Sama seperti versi model kredit dicatat pada tiap rekomendasi di demo.'\}
\item \{'h': 'Jalankan kumpulan uji sebelum berpindah versi', 'p': 'Perpindahan versi adalah perubahan perangkat lunak, dan pantas diperlakukan seperti perubahan perangkat lunak.'\}
\end{isteps}
\begin{iband}[signal]
Ini kelas kegagalan yang tidak muncul di pengembangan dan hanya muncul di produksi --- dan gejalanya \textbf{selalu} terlihat seperti ``tiba-tiba agennya jadi bodoh''.
\end{iband}
\end{frame}

\framekicker{Operasional}
\begin{frame}[shrink=25]{Batas laju: yang membatasi bukan kecepatan, melainkan token per menit}
{\fontsize{9.4}{12.6}\selectfont\color{fg2}Penyedia membatasi dua hal sekaligus: permintaan per menit dan \textbf{token per menit}. Agen dengan konteks panjang menabrak yang kedua jauh sebelum yang pertama.\par}\vskip 4pt
\begin{minipage}[t]{0.485\textwidth}
{\fontsize{9.4}{12.6}\selectfont\color{fg2}\textbf{Kenapa ini mengejutkan orang}\par}\vskip 4pt
{\fontsize{9.4}{12.6}\selectfont\color{fg2}Sepuluh pengguna bersamaan terdengar kecil. Sepuluh pengguna $\times$ enam giliran $\times$ konteks yang menumpuk bisa berarti ratusan ribu token per menit.\par}\vskip 4pt

\end{minipage}%
\hfill
\begin{minipage}[t]{0.485\textwidth}
{\fontsize{9.4}{12.6}\selectfont\color{fg2}\textbf{Yang harus ada sebelum dipakai banyak orang}\par}\vskip 4pt
\begin{ibullets}
\item Antrean, bukan pengulangan langsung
\item Mundur bertahap ketika ditolak
\item Batas jumlah proses agen yang berjalan bersamaan
\end{ibullets}

\end{minipage}%
\begin{iband}[signal]
Dan satu akibat desain: \textbf{memperpendek konteks menaikkan kapasitas}, bukan hanya menurunkan biaya. Keduanya obat yang sama.
\end{iband}
\end{frame}

\framekicker{Operasional}
\begin{frame}[shrink=25]{Model di mesin sendiri: kapan masuk akal}
\begin{minipage}[t]{0.494\textwidth}
\begin{icardaccent}
\cardh{Ketika datanya yang menentukan}
Ada data yang tidak boleh keluar dari perimeter. Ini keputusan \textbf{kepatuhan}, dan ia mengalahkan pertimbangan biaya maupun kualitas.
\end{icardaccent}
\end{minipage}%
\hfill
\begin{minipage}[t]{0.494\textwidth}
\begin{icardbad}
\cardh{Bukan karena lebih murah}
Sering justru tidak, kalau pemakaiannya tidak rata: GPU menganggur tetap dibayar, sedangkan API yang tidak dipanggil tidak.
\end{icardbad}
\end{minipage}%
{\fontsize{9.4}{12.6}\selectfont\color{fg2}Yang berubah secara teknis: model terbuka umumnya lebih lemah pada \textbf{kepatuhan skema alat} dan lebih mudah berhenti terlalu awal. Itu bukan alasan menolaknya --- itu alasan mengukurnya pada kumpulan uji Anda sebelum memutuskan.\par}\vskip 4pt
\begin{iband}[signal]
Dijalankan sungguhan di demo dengan \inlinecode{qwen3:8b}: gelungnya benar, tetapi \textbf{berhenti satu langkah lebih awal} tanpa memanggil alat pengiriman rekomendasi. Kelemahan yang terukur, bukan yang ditebak.
\end{iband}
\end{frame}

\framekicker{Operasional}
\begin{frame}[shrink=25]{Enam angka yang pantas dicatat pada tiap proses}
\begin{itable}{>{\raggedright\arraybackslash}p{0.2256\textwidth}>{\raggedright\arraybackslash}p{0.3572\textwidth}>{\raggedright\arraybackslash}p{0.3572\textwidth}}
\thead{Angka} & \thead{Yang diberitahukannya} & \thead{Gejala kalau memburuk} \\ \midrule
Token masukan per tugas & Biaya sebenarnya & Naik pelan tanpa perubahan kode $\rightarrow$ riwayat tidak dipangkas \\
\inlinecode{cache\_read} bukan nol & Singgahan menyala & Jadi nol $\rightarrow$ ada yang menaruh sesuatu yang berubah di awalan \\
Giliran per tugas & Panjang penalaran & Naik $\rightarrow$ alat berubah, atau deskripsinya jadi kabur \\
Laju penolakan skema & Kecocokan model & Naik sesudah ganti model $\rightarrow$ modelnya, bukan agennya \\
Waktu alat lawan waktu model & Ke mana waktu pergi & Waktu alat mendominasi $\rightarrow$ optimasi model tidak akan terasa \\
Laju eskalasi ke manusia & Kejujuran sistem & Turun ke nol $\rightarrow$ curigai, jangan rayakan \\
\end{itable}
{\fontsize{9.4}{12.6}\selectfont\color{fg2}Tidak satu pun dari enam ini menimbulkan \textbf{galat}. Semuanya bergeser diam-diam, dan hanya terlihat kalau dicatat sejak awal.\par}\vskip 4pt
\end{frame}

\framekicker{Penutup}
\begin{frame}[shrink=25]{Yang dibawa pulang dari bab ini}
\begin{isteps}
\item \{'h': 'Token adalah satuan segalanya', 'p': 'Tagihan, waktu tunggu, dan batas konteks --- ketiganya diukur dengan satuan yang sama, dan agen adalah beban \textbf{masukan}.'\}
\item \{'h': 'Percakapan sepuluh giliran menagih enam kali isinya', 'p': 'Sebab tiap giliran membayar ulang semua giliran sebelumnya. Mengurangi langkah mengalahkan memangkas kata.'\}
\item \{'h': 'Model tidak mengingat; daftar pesan Anda yang mengingat', 'p': 'Itu kabar baik --- keadaan yang Anda pegang bisa diperiksa dan dipotong dengan sengaja.'\}
\item \{'h': 'Panggilan alat adalah niat, bukan eksekusi', 'p': 'Semua batas izin hidup di celah antara model menuliskannya dan kode Anda menjalankannya.'\}
\item \{'h': 'Pilih model dengan kumpulan uji Anda, bukan papan skor', 'p': 'Kepatuhan skema, kemauan berhenti, dan kepatuhan instruksi panjang tidak ada di sana.'\}
\end{isteps}
\end{frame}
\note{Kalau waktunya mepet, tiga slide yang tidak boleh dilewat: pertumbuhan konteks, panggilan alat sebagai keluaran terstruktur, dan tiga sifat model. Sisanya bisa jadi bacaan.}

\end{document}
